This thesis presented a complete, working chain from force sensing to traction output. Tensile and compressive loads were captured with bonded electrical resistance strain gauges in a Wheatstone configuration, conditioned by a HX711 module (PGA + 24-bit ADC), filtered and shaped in software on the Hypatia board using a moving average filter and PD control, converted to analog by a MCP4725 \ac(DAC), amplified with an LM358 to a clean ±10 V command, and finally applied to a discrete transistor H-bridge that drives the DC motors. A phone-based dashboard enabled live visualization and parameter tuning, and wheel encoders provided feedback for monitoring and future closed-loop refinement. In practice, the demonstrator translates measured load into stable, bidirectional wheel torque and exposes each processing stage for inspection.

All primary requirements were met. Load sensing and basic calibration were achieved. Signal conditioning and control (filter + PD) reduced noise and improved responsiveness. Actuation from \ac{DAC} to amplifier to H-bridge operated as intended. Feedback from both wheel encoders was acquired and system integration emphasized usability with separated power domains, a master switch, status indication, and a mobile dashboard. The demonstrator therefore satisfies the functional and integration goals defined at the outset.

A brief note on accuracy and data loss: residual uncertainty comes mainly from sensor and bridge tolerances (bonding, mismatch, temperature drift), electrical limits (ADC/DAC resolution, reference and op-amp behavior) and  drive interactions (H-bridge dead zones, motor non-linearities). The implemented filtering and PD action mitigate several of these effects, but they set a practical floor for precision.

\subsection{Next Steps}
The next development step is a more systematic evaluation of the acquired signals and a tighter alignment between measurement, control, and actuation. Beyond extended testing, several practical refinements will raise repeatability and ease of use: joystick calibration (center, deadband, linearity), trim potentiometers on the strain-gauge bridge for fine span adjustment, and a small potentiometer in the voltage divider of the \ac{DAC} path to set the analogue scaling precisely.

The main architectural change concerns the drive stage. Although the LM358 currently delivers a clean $\pm 10~\text{V}$ control signal, the discrete H-bridge that follows effectively provides 12 V with reversible polarity but does not implement a proportional voltage (or torque) response to that analogue command. In other words, direction is achieved, but the desired change in motor voltage (and thus speed/torque) does not yet track the LM358 output.

One idea to resolve this, the discrete H-bridge could be replaced by a motor driver for example an L298N and the command translation moves into software. The LM358 output (or the upstream digital value) can be mapped to a \ac{PWM} duty cycle and a direction bit. The motor driver then applies the required current from the 12V rail with proper switching, protection, and current capability. 

In summary, the immediate focus is deeper testing of the raw/filtered/PD signals, calibration aids on the sensor and \ac{DAC} scaling, as well as replacing the discrete H-bridge with a PWM-driven motor driver so that the motor voltage—and thus the demonstrator’s behavior,accurately follows the intended command.
\newpage